\section{Funcionalidades Detalhadas}

\subsection{Metodologia de Organizacao Financeira}

\subsubsection{Modelos de Organizacao Suportados}
O sistema oferece flexibilidade para diferentes tipos de organizacao financeira:

\paragraph{Modelo 1: Renda Conjunta (Padrao)}
\begin{itemize}
    \item \textbf{Aplicacao}: Casais que juntam todas as rendas
    \item \textbf{Divisao}: 50/50 apos todas as alocacoes
    \item \textbf{Exemplo}: Casal soma R\$ 8.000, apos gastos fixos, poupanca etc., sobram R\$ 2.000 que sao divididos em R\$ 1.000 para cada
\end{itemize}

\paragraph{Modelo 2: Divisao Igualitaria com Acerto (Splitwise)}
\begin{itemize}
    \item \textbf{Aplicacao}: Casal que quer dividir gastos compartilhados de forma justa
    \item \textbf{Metodo}: Cada pessoa cadastra o que pagou, sistema calcula divisao 50/50
    \item \textbf{Exemplo}: Joao gastou R\$ 1.000, Maria R\$ 400 -> Maria deve R\$ 300 para Joao
    \item \textbf{Acerto}: Sistema sugere transferencia mensal para equilibrar as contas
\end{itemize}

\paragraph{Modelo 3: Organizacao Individual}
\begin{itemize}
    \item \textbf{Aplicacao}: Republicas, pessoas solteiras, ou independencia financeira total
    \item \textbf{Divisao}: Cada pessoa responsavel por suas proprias contas
    \item \textbf{Exemplo}: Divisao apenas de gastos compartilhados (aluguel, internet)
\end{itemize}

\paragraph{Modelo 4: Contas Separadas + Fundo Comum}
\begin{itemize}
    \item \textbf{Aplicacao}: Casais que querem autonomia financeira mas compartilhar despesas essenciais
    \item \textbf{Divisao}: Cada um contribui valor fixo para fundo comum, restante fica individual
    \item \textbf{Exemplo}: Casal contribui R\$ 2.000 cada para fundo comum, mantem resto individual
    \item \textbf{Vantagem}: Transparencia apenas no essencial, total privacidade no resto
\end{itemize}

\subsubsection{Fluxo de Organizacao Adaptavel}
O sistema implementa a seguinte metodologia que se adapta ao modelo escolhido:

\begin{enumerate}
    \item \textbf{Receitas}: Individuais ou conjuntas conforme modelo
    \item \textbf{Gastos Fixos}: Pessoais e/ou compartilhados
    \item \textbf{Custos Variaveis}: Orcamento individual ou conjunto
    \item \textbf{Poupanca}: Metas individuais ou conjuntas
    \item \textbf{Reserva de Emergencia}: Por pessoa ou familia
    \item \textbf{Investimentos}: Estrategia individual ou conjunta
    \item \textbf{Gastos Livres}: Calculo adaptado ao modelo
\end{enumerate}

\subsubsection{Implementacao Tecnica}

\begin{lstlisting}[language=C, caption=Motor de Calculo Adaptativo]
const calcularOrganizacaoFinanceira = (receitas, gastos, modelo, configuracao) => {
  let resultado = {};
  
  switch(modelo) {
    case 'RENDA_CONJUNTA':
      resultado = calcularRendaConjunta(receitas, configuracao);
      break;
    case 'DIVISAO_SPLITWISE':
      resultado = calcularDivisaoSplitwise(gastos, configuracao);
      break;
    case 'INDIVIDUAL':
      resultado = calcularIndividual(receitas, configuracao);
      break;
    case 'FUNDO_COMUM':
      resultado = calcularFundoComum(receitas, configuracao);
      break;
  }
  
  return resultado;
};

const calcularRendaConjunta = (receitas, config) => {
  const totalReceitas = receitas.reduce((sum, r) => sum + r.valor, 0);
  const gastosFixes = calcularGastosFixos();
  const custosVariaveis = config.percentualVariaveis * totalReceitas;
  const poupanca = config.percentualPoupanca * totalReceitas;
  const reservaEmergencia = config.percentualReserva * totalReceitas;
  const investimentos = config.percentualInvestimentos * totalReceitas;
  
  const valorUtilizado = gastosFixes + custosVariaveis + poupanca + 
                        reservaEmergencia + investimentos;
  const gastosLivres = totalReceitas - valorUtilizado;
  const gastosLivresIndividual = gastosLivres / 2; // Divisao 50/50
  
  return { gastosLivresIndividual, totalReceitas, ... };
};

const calcularDivisaoSplitwise = (gastos, config) => {
  // Gastos registrados por cada pessoa
  const gastosPessoa1 = gastos.filter(g => g.pagador === 'pessoa1');
  const gastosPessoa2 = gastos.filter(g => g.pagador === 'pessoa2');
  
  // Total gasto por cada pessoa
  const totalPessoa1 = gastosPessoa1.reduce((sum, g) => sum + g.valor, 0);
  const totalPessoa2 = gastosPessoa2.reduce((sum, g) => sum + g.valor, 0);
  
  // Total geral e divisao ideal (50/50)
  const totalGeral = totalPessoa1 + totalPessoa2;
  const idealPorPessoa = totalGeral / 2;
  
  // Calculo de quem deve o que
  const saldoPessoa1 = idealPorPessoa - totalPessoa1;
  const saldoPessoa2 = idealPorPessoa - totalPessoa2;
  
  // Quem deve transferir para quem
  let transferencia = null;
  if (saldoPessoa1 > 0) {
    transferencia = {
      de: 'pessoa2',
      para: 'pessoa1', 
      valor: saldoPessoa1
    };
  } else if (saldoPessoa2 > 0) {
    transferencia = {
      de: 'pessoa1',
      para: 'pessoa2',
      valor: saldoPessoa2
    };
  }
  
  return { 
    totalGeral,
    idealPorPessoa,
    totalPessoa1,
    totalPessoa2,
    transferencia
  };
};
        \item Maria pagou: Energia R\$ 200, Internet R\$ 100, Uber Eats R\$ 80
    \end{itemize}
    \item \textbf{Calculo automatico de divisao justa}:
    \begin{itemize}
        \item Total gasto: R\$ 2.210 (R\$ 1.830 Joao + R\$ 380 Maria)
        \item Ideal por pessoa: R\$ 1.105 cada
        \item Maria deve para Joao: R\$ 725 (R\$ 1.105 - R\$ 380)
    \end{itemize}
    \item \textbf{Funcionalidades Splitwise}:
    \begin{itemize}
        \item Historico de quem deve o que
        \item Sugestao de acerto mensal
        \item Categorias de gastos compartilhados vs individuais
        \item Registro de pagamentos para equilibrar contas
    \end{itemize}
\end{itemize}
O sistema oferece flexibilidade para diferentes tipos de organizacao financeira:

\paragraph{Modelo 1: Renda Conjunta (Padrao)}
\begin{itemize}
    \item \textbf{Aplicacao}: Casais que juntam todas as rendas
    \item \textbf{Divisao}: 50/50 apos todas as alocacoes
    \item \textbf{Exemplo}: Casal soma R\$ 8.000, apos gastos fixos, poupanca etc., sobram R\$ 2.000 que sao divididos em R\$ 1.000 para cada
\end{itemize}

\paragraph{Modelo 2: Divisao Igualitaria com Acerto (Splitwise)}
\begin{itemize}
    \item \textbf{Aplicacao}: Casal que quer dividir gastos compartilhados de forma justa
    \item \textbf{Metodo}: Cada pessoa cadastra o que pagou, sistema calcula divisao 50/50
    \item \textbf{Exemplo}: Joao gastou R\$ 1.000, Maria R\$ 400 -> Maria deve R\$ 300 para Joao
    \item \textbf{Acerto}: Sistema sugere transferencia mensal para equilibrar as contas
\end{itemize}

\paragraph{Modelo 3: Organizacao Individual}
\begin{itemize}
    \item \textbf{Aplicacao}: Republicas, pessoas solteiras, ou independencia financeira total
    \item \textbf{Divisao}: Cada pessoa responsavel por suas proprias contas
    \item \textbf{Exemplo}: Divisao apenas de gastos compartilhados (aluguel, internet)
\end{itemize}

\subsubsection{Fluxo de Organizacao Adaptavel}
O sistema implementa a seguinte metodologia que se adapta ao modelo escolhido:

\begin{enumerate}
    \item \textbf{Receitas}: Individuais ou conjuntas conforme modelo
    \item \textbf{Gastos Fixos}: Pessoais e/ou compartilhados
    \item \textbf{Custos Variaveis}: Orcamento individual ou conjunto
    \item \textbf{Poupanca}: Metas individuais ou conjuntas
    \item \textbf{Reserva de Emergencia}: Por pessoa ou familia
    \item \textbf{Investimentos}: Estrategia individual ou conjunta
    \item \textbf{Gastos Livres}: Calculo adaptado ao modelo
\end{enumerate}

\subsubsection{Calculo Automatico Adaptavel}
\begin{lstlisting}[caption=Exemplo de calculo automatico por modelo]
// Logica de calculo adaptavel por modelo de organizacao
const calcularDistribuicaoFinanceira = (receitas, configuracao, modelo) => {
  let resultado = {};
  
  switch(modelo) {
    case 'RENDA_CONJUNTA':
      resultado = calcularRendaConjunta(receitas, configuracao);
      break;
    case 'DIVISAO_SPLITWISE':
      resultado = calcularDivisaoSplitwise(gastos, configuracao);
      break;
    case 'INDIVIDUAL':
      resultado = calcularIndividual(receitas, configuracao);
      break;
    case 'FUNDO_COMUM':
      resultado = calcularFundoComum(receitas, configuracao);
      break;
  }
  
  return resultado;
};

const calcularRendaConjunta = (receitas, config) => {
  const totalReceitas = receitas.reduce((sum, r) => sum + r.valor, 0);
  const gastosFixes = calcularGastosFixos();
  const custosVariaveis = config.percentualVariaveis * totalReceitas;
  const poupanca = config.percentualPoupanca * totalReceitas;
  const reservaEmergencia = config.percentualReserva * totalReceitas;
  const investimentos = config.percentualInvestimentos * totalReceitas;
  
  const valorUtilizado = gastosFixes + custosVariaveis + poupanca + 
                        reservaEmergencia + investimentos;
  const gastosLivres = totalReceitas - valorUtilizado;
  const gastosLivresIndividual = gastosLivres / 2; // Divisao 50/50
  
  return { gastosLivresIndividual, totalReceitas, ... };
};

const calcularDivisaoSplitwise = (gastos, config) => {
  // Gastos registrados por cada pessoa
  const gastosPessoa1 = gastos.filter(g => g.pagador === 'pessoa1');
  const gastosPessoa2 = gastos.filter(g => g.pagador === 'pessoa2');
  
  // Total gasto por cada pessoa
  const totalPessoa1 = gastosPessoa1.reduce((sum, g) => sum + g.valor, 0);
  const totalPessoa2 = gastosPessoa2.reduce((sum, g) => sum + g.valor, 0);
  
  // Total geral e divisao ideal (50/50)
  const totalGeral = totalPessoa1 + totalPessoa2;
  const idealPorPessoa = totalGeral / 2;
  
  // Calculo de quem deve o que
  const saldoPessoa1 = idealPorPessoa - totalPessoa1;
  const saldoPessoa2 = idealPorPessoa - totalPessoa2;
  
  // Quem deve transferir para quem
  let transferencia = null;
  if (saldoPessoa1 > 0) {
    transferencia = {
      de: 'pessoa2',
      para: 'pessoa1', 
      valor: saldoPessoa1
    };
  } else if (saldoPessoa2 > 0) {
    transferencia = {
      de: 'pessoa1',
      para: 'pessoa2',
      valor: saldoPessoa2
    };
  }
  
  return { 
    totalGeral,
    idealPorPessoa,
    totalPessoa1,
    totalPessoa2,
    transferencia
  };
};

const calcularFundoComum = (receitas, config) => {
  const fundoComumPorPessoa = config.valorFundoComum;
  const totalFundoComum = fundoComumPorPessoa * 2;
  
  // Cada pessoa contribui valor fixo para o fundo
  const resultadoPessoa1 = {
    receitaTotal: receitas.pessoa1.total,
    contribuicaoFundo: fundoComumPorPessoa,
    dinheiroIndividual: receitas.pessoa1.total - fundoComumPorPessoa
  };
  
  const resultadoPessoa2 = {
    receitaTotal: receitas.pessoa2.total,
    contribuicaoFundo: fundoComumPorPessoa,
    dinheiroIndividual: receitas.pessoa2.total - fundoComumPorPessoa
  };
  
  return { 
    fundoComum: totalFundoComum,
    pessoa1: resultadoPessoa1,
    pessoa2: resultadoPessoa2 
  };
};
\end{lstlisting}

\subsubsection{Modulo de Configuracao de Modelo}

\paragraph{Selecao do Modelo de Organizacao}
\begin{itemize}
    \item \textbf{Onboarding Inteligente}:
    \begin{itemize}
        \item Questionario inicial sobre situacao familiar/habitacional
        \item Sugestao automatica do modelo mais adequado
        \item Possibilidade de mudanca posterior
    \end{itemize}
    
    \item \textbf{Configuracoes por Modelo}:
    \begin{itemize}
        \item \textit{Renda Conjunta}: Percentuais de alocacao, divisao 50/50 ou customizada
        \item \textit{Divisao Splitwise}: Categorias de gastos compartilhados, configuracao de acerto mensal
        \item \textit{Individual}: Apenas gastos compartilhados configuraveis
        \item \textit{Fundo Comum}: Valor da contribuicao mensal por pessoa
    \end{itemize}
    
    \item \textbf{Flexibilidade de Mudanca}:
    \begin{itemize}
        \item Migracao entre modelos sem perda de dados historicos
        \item Simulacao de cenarios antes da mudanca
        \item Backup automatico das configuracoes anteriores
    \end{itemize}
\end{itemize}

\paragraph{Casos de Uso Detalhados}

\textbf{Cenario 1: Casal com Renda Conjunta (Voce e sua esposa)}
\begin{itemize}
    \item Soma total: R\$ 8.000 (R\$ 5.000 + R\$ 3.000)
    \item Gastos fixos: R\$ 3.500 (aluguel, contas, etc.)
    \item Poupanca (15\%): R\$ 1.200
    \item Reserva (10\%): R\$ 800
    \item Investimentos (15\%): R\$ 1.200
    \item Gastos livres restantes: R\$ 1.300 -> R\$ 650 para cada
\end{itemize}

\textbf{Cenario 2: Casal com Divisao Proporcional}
\begin{itemize}
    \item Pessoa A: R\$ 5.000 (62.5\% da renda total)
    \item Pessoa B: R\$ 3.000 (37.5\% da renda total)
    \item Aluguel R\$ 1.500: A paga R\$ 937, B paga R\$ 563
    \item Energia R\$ 200: A paga R\$ 125, B paga R\$ 75
    \item Cada um gerencia o restante individualmente
\end{itemize}

\textbf{Cenario 3: Republica/Individual}
\begin{itemize}
    \item 3 pessoas dividindo apartamento
    \item Gastos compartilhados: aluguel, internet, limpeza
    \item Divisao igualitaria ou por uso/quarto
    \item Cada pessoa gerencia 100\% de suas financas pessoais
\end{itemize}

\subsubsection{Modulo de Receitas}

\paragraph{Funcionalidades}
\begin{itemize}
    \item \textbf{Cadastro de Receitas}:
    \begin{itemize}
        \item Descricao personalizada
        \item Valor em reais (R\$)
        \item Categoria (salario, freelance, investimentos, etc.)
        \item Data de recebimento
        \item Recorrencia (mensal, eventual)
        \item Status (recebida, pendente, cancelada)
    \end{itemize}
    
    \item \textbf{Categorias Pre-definidas}:
    \begin{itemize}
        \item Salario CLT
        \item Freelance/Autonomo
        \item Rendimentos de Investimentos
        \item Alugueis Recebidos
        \item Vendas/Comissoes
        \item Outras Receitas
    \end{itemize}
    
    \item \textbf{Relatorios de Receitas}:
    \begin{itemize}
        \item Evolucao mensal
        \item Distribuicao por categoria
        \item Previsao vs realizado
        \item Sazonalidade
    \end{itemize}
\end{itemize}

\subsubsection{Modulo de Gastos Fixos}

\paragraph{Funcionalidades}
\begin{itemize}
    \item \textbf{Cadastro de Gastos Fixos}:
    \begin{itemize}
        \item Descricao (ex: Aluguel, Internet, Seguro)
        \item Valor mensal
        \item Categoria
        \item Data de vencimento
        \item Conta de debito
        \item Recorrencia automatica
    \end{itemize}
    
    \item \textbf{Categorias de Gastos Fixos}:
    \begin{itemize}
        \item Moradia (aluguel, condominio, IPTU)
        \item Transporte (financiamento, seguro, IPVA)
        \item Alimentacao (supermercado basico)
        \item Saude (plano de saude, medicamentos)
        \item Educacao (mensalidades, cursos)
        \item Servicos (internet, telefone, streaming)
        \item Seguros (vida, residencial, veicular)
    \end{itemize}
    
    \item \textbf{Lembretes e Alertas}:
    \begin{itemize}
        \item Notificacoes de vencimento
        \item Alertas de aumento de valores
        \item Sugestoes de economia
    \end{itemize}
\end{itemize}

\subsubsection{Modulo de Custos Variaveis}

\paragraph{Funcionalidades}
\begin{itemize}
    \item \textbf{Orcamento por Categoria}:
    \begin{itemize}
        \item Definicao de limite mensal
        \item Acompanhamento em tempo real
        \item Alertas de aproximacao do limite
    \end{itemize}
    
    \item \textbf{Categorias de Custos Variaveis}:
    \begin{itemize}
        \item Alimentacao (restaurantes, delivery)
        \item Transporte (combustivel, Uber, transporte publico)
        \item Lazer (cinema, shows, viagens)
        \item Vestuario
        \item Saude (consultas particulares, farmacia)
        \item Casa (manutencao, decoracao)
        \item Educacao (livros, cursos extras)
    \end{itemize}
    
    \item \textbf{Controle de Gastos}:
    \begin{itemize}
        \item Registro rapido via app
        \item Foto do comprovante
        \item Divisao por tags
        \item Geolocalizacao opcional
    \end{itemize}
\end{itemize}

\subsubsection{Modulo de Poupanca e Reservas}

\paragraph{Conceito e Diferenciacao}
O sistema oferece **indicacoes percentuais** para alocacao financeira, funcionando como um organizador digital que **substitui planilhas Excel**:

\begin{itemize}
    \item \textbf{Poupanca (Meta/Objetivo)}: Recursos destinados para objetivos especificos de medio/longo prazo
    \item \textbf{Reserva de Emergencia}: Valor separado para situacoes imprevistas e urgencias
    \item \textbf{Diferenca Principal}: Poupanca tem objetivo especifico; Reserva e para emergencias
\end{itemize}

\paragraph{Funcionalidades}
\begin{itemize}
    \item \textbf{Indicacao Percentual para Poupanca}:
    \begin{itemize}
        \item Sugestao de 15-20\% da renda liquida
        \item Objetivos especificos (viagem, casa, carro, casamento)
        \item Valor total da meta definido pelo usuario
        \item Prazo estimado para atingir o objetivo
        \item Calculo automatico: se meta = R\$ 30.000 e poupa R\$ 500/mes = 60 meses
        \item Progresso visual em barras e percentuais
    \end{itemize}
    
    \item \textbf{Indicacao Percentual para Reserva de Emergencia}:
    \begin{itemize}
        \item Sugestao de 10-15\% da renda liquida
        \item Meta recomendada: 6-12 meses de gastos essenciais
        \item Calculo automatico: se gastos essenciais = R\$ 3.000, reserva = R\$ 18.000-36.000
        \item **Sem acompanhamento de rendimento** - apenas controle de valores salvos
        \item Simulacao de tempo para atingir meta com valor mensal definido
    \end{itemize}
    
    \item \textbf{Controle Manual de Valores}:
    \begin{itemize}
        \item Registro manual de depositos realizados
        \item Registro manual de retiradas (com justificativa)
        \item Saldo atual informado pelo usuario
        \item Historico de movimentacoes para acompanhamento
    \end{itemize}
\end{itemize}

\subsubsection{Modulo de Investimentos}

\paragraph{Conceito e Diferenciacao}
O modulo de investimentos funciona como um **indicador percentual** e organizador de valores, **sem integracao com corretoras ou APIs**:

\begin{itemize}
    \item \textbf{Diferenca da Poupanca}: Investimentos buscam crescimento do capital; Poupanca e para objetivos especificos
    \item \textbf{Diferenca da Reserva}: Investimentos podem ter risco; Reserva deve ser conservadora e liquida
    \item \textbf{Funcao no FinOps}: Sugestao de quanto alocar e organizacao manual dos valores
\end{itemize}

\paragraph{Funcionalidades}
\begin{itemize}
    \item \textbf{Indicacao Percentual para Investimentos}:
    \begin{itemize}
        \item Sugestao de 10-25\% da renda liquida (dependendo do perfil)
        \item Educacao sobre perfil conservador (10-15\%), moderado (15-20\%), arrojado (20-25\%)
        \item Calculo automatico do valor mensal disponivel para investir
    \end{itemize}
    
    \item \textbf{Categorizacao de Tipos (Educacional)}:
    \begin{itemize}
        \item \textbf{Renda Fixa}: CDB, LCI, LCA, Tesouro Direto (menor risco)
        \item \textbf{Renda Variavel}: Acoes, FIIs, ETFs (maior risco e potencial)
        \item \textbf{Fundos}: Fundos de Investimento diversos
        \item \textbf{Previdencia}: PGBL, VGBL (longo prazo)
        \item \textbf{Outros}: Criptomoedas, commodities
        \item **Apenas para organizacao e educacao - sem acompanhamento automatico**
    \end{itemize}
    
    \item \textbf{Controle Manual de Aportes}:
    \begin{itemize}
        \item Registro manual de valores investidos por mes
        \item Registro por categoria de investimento
        \item Total acumulado informado pelo usuario
        \item **Sem acompanhamento de rentabilidade automatica**
        \item Historico de aportes para controle pessoal
    \end{itemize}
    
    \item \textbf{Educacao Financeira}:
    \begin{itemize}
        \item Explicacao basica de cada tipo de investimento
        \item Sugestao de diversificacao por perfil
        \item Calculadoras simples de juros compostos
        \item Dicas de quando revisar e rebalancear (manualmente)
    \end{itemize}
\end{itemize}

\subsubsection{Modulo de Gastos Livres}

\paragraph{Funcionalidades}
\begin{itemize}
    \item \textbf{Calculo Automatico}:
    \begin{itemize}
        \item Valor disponivel apos todas as alocacoes
        \item Divisao automatica para casal (50/50)
        \item Ajuste manual se necessario
    \end{itemize}
    
    \item \textbf{Controle Individual}:
    \begin{itemize}
        \item Saldo individual em tempo real
        \item Historico de gastos pessoais
        \item Categorizacao livre
        \item Transferencias entre parceiros
    \end{itemize}
    
    \item \textbf{Flexibilidade}:
    \begin{itemize}
        \item Acumulo para meses seguintes
        \item Antecipacao de valores futuros
        \item Emprestimos internos entre parceiros
    \end{itemize}
\end{itemize}

\subsection{Funcionalidades Transversais}

\subsubsection{Dashboard e Relatorios}
\begin{itemize}
    \item Visao geral financeira em tempo real
    \item Graficos de evolucao temporal
    \item Comparacao entre periodos
    \item Indicadores de performance financeira
    \item Export para PDF e Excel
\end{itemize}

\subsubsection{Notificacoes}
\begin{itemize}
    \item Push notifications para mobile
    \item Email notifications opciais
    \item Lembretes de vencimentos
    \item Alertas de orcamento
    \item Resumos semanais/mensais
\end{itemize}

\subsubsection{Configuracoes}
\begin{itemize}
    \item Personalizacao de categorias
    \item Definicao de percentuais padrao
    \item Configuracao de lembretes
    \item Preferencias de relatorios
    \item Configuracoes de privacidade
\end{itemize}
